\section{Sources et reproduction}

\subsection{D'où viennent les chiffres}

\begin{itemize}[leftmargin=1.2em,itemsep=3pt,topsep=2pt]
  \item \textbf{Checkpoints supervisés} (Lightning) : les champs
        \code{epoch}, \code{global\_step} et l'état des boucles
        (\code{loops.fit\_loop}) donnent le nombre exact de steps, le nombre
        de lots par époque et les compteurs d'époque. Lecture par
        \code{torch.load} avec \code{weights\_only=False}.
  \item \textbf{Checkpoints self-play} (\code{.pt}) : les clés
        \code{iteration} et \code{global\_step} donnent l'état exact de chaque
        jalon (par exemple \code{iteration=436}, \code{global\_step=26356}).
  \item \textbf{Shards binaires} : les tailles et les tenseurs lus dans
        \chemin{Downloads/dataset\_shards} donnent la taille d'époque du
        prétraining (chaque shard de 500 parties contient environ 16\,600
        positions).
  \item \textbf{Wandb} : le projet
        \href{https://wandb.ai/leclercqmathieu14/alphazero-chess}{%
        leclercqmathieu14/alphazero-chess} contient les 68 runs de mars et
        avril. Un export local complet a été réalisé dans
        \chemin{out/wandb\_export} (configurations, résumés, historiques CSV,
        journaux de sortie).
\end{itemize}

\subsection{Exporter wandb en local}

L'export a été produit par \chemin{out/fetch\_wandb.py} (clé API lue depuis
\chemin{out/wandb\_key.txt}, ce dossier est ignoré par git), puis agrégé par
\chemin{out/aggregate.py} en deux fichiers de travail :
\chemin{merged\_selfplay.csv} (une ligne par itération) et
\chemin{runs\_augmented.json} (une entrée par lancement, avec la plage
d'itérations couverte).

\subsection{Regénérer les figures}

Depuis le dossier \chemin{docs/rapports/2026-04-entrainement-iter436} :

\begin{verbatim}
python figures/generer_figures.py
\end{verbatim}

Le script lit \chemin{out/wandb\_export} et écrit les neuf figures PDF dans
\chemin{figures/}. Il demande \code{matplotlib} et \code{numpy}
(\code{python3} sur la machine de rédaction).

\subsection{Documentation liée}

\begin{itemize}[leftmargin=1.2em,itemsep=2pt,topsep=2pt]
  \item \chemin{docs/training-history.md} : la même reconstitution en
        anglais, au format court.
  \item \chemin{python\_src/train\_supervised.py} et
        \chemin{python\_src/train\_self\_play.py} : les deux scripts
        d'entraînement.
  \item \chemin{python\_src/checkpoints/} : les jalons conservés de la
        génération 1 et de la génération 2.
  \item \chemin{CHANGELOG.md} : les jalons techniques du projet depuis 2022.
\end{itemize}
